\let\origsectionmark\sectionmark\renewcommand*{\sectionmark}[1]{\origsectionmark{Absorption by current and spin current}}% short running head
\section[Absorption by the current and spin current of a superconductor]{Absorption by the uniform current and spin current of a superconductor}\label{app:D.40}
\let\sectionmark\origsectionmark

A boson absorbed at $\vect{q} \to 0$ couples to a superconducting metal through the total velocity of its electrons, if it is a dark photon or a photon, and through their spin current, if it is an axion-like particle. This section derives what \cref{sec:15.4} uses: the Nambu form of both operators, the absorption by the spin current of a clean superconductor, the exact second-order amplitude of phonon-assisted absorption and its conductivity, the limit of a Debye transport function, and impurity-assisted absorption in the dirty limit. The notation is that of \cref{sec:7.2,app:D.20}: the band energy $\xi_{\vect{k}}$ from the chemical potential, the gap $\Delta_{\mathrm{BCS}}$, the quasiparticle energy $E_{\vect{k}}$ and coherence factors $u^{\mathrm{B}}_{\vect{k}}$, $v^{\mathrm{B}}_{\vect{k}}$ of \cref{eq:7.2-E}, the quasiparticle operators $\gamma_{\vect{k}1}$, $\gamma_{\vect{k}2}$ of \cref{eq:7.2-gamma}, and the Nambu spinor $\Psi_{\vect{k}} = (c_{\vect{k}\uparrow}, c^\dagger_{-\vect{k}\downarrow})^{\mathrm{T}}$ with the matrices of \cref{eq:7.2-nambu}. Throughout, $T = 0$, one band crosses the Fermi level, and the crystal is symmetric under time reversal, so that $\xi_{-\vect{k}} = \xi_{\vect{k}}$ and the band velocity $\vect{v}_{\vect{k}} = \partial\xi_{\vect{k}}/\partial\vect{k}$ obeys $\vect{v}_{-\vect{k}} = -\vect{v}_{\vect{k}}$.

\subsection*{Uniform one-body operators in Nambu form}

Take a uniform operator that is diagonal in $\vect{k}$ and in spin, with a weight $f_{\vect{k}}$ that is odd, $f_{-\vect{k}} = -f_{\vect{k}}$, and a spin factor $s_\sigma$: $O = \sum_{\vect{k}\sigma}s_\sigma f_{\vect{k}}\,c^\dagger_{\vect{k}\sigma}c_{\vect{k}\sigma}$. In the spin-down term relabel $\vect{k} \to -\vect{k}$, so that $\sum_{\vect{k}}f_{\vect{k}}\,c^\dagger_{\vect{k}\downarrow}c_{\vect{k}\downarrow} = -\sum_{\vect{k}}f_{\vect{k}}\,c^\dagger_{-\vect{k}\downarrow}c_{-\vect{k}\downarrow}$. With $c^\dagger_{-\vect{k}\downarrow}c_{-\vect{k}\downarrow} = 1 - c_{-\vect{k}\downarrow}c^\dagger_{-\vect{k}\downarrow}$, and since $\sum_{\vect{k}}f_{\vect{k}} = 0$,
\begin{equation}
  O = \sum_{\vect{k}}f_{\vect{k}}\bigl[s_\uparrow c^\dagger_{\vect{k}\uparrow}c_{\vect{k}\uparrow} + s_\downarrow c_{-\vect{k}\downarrow}c^\dagger_{-\vect{k}\downarrow}\bigr] .
  \label{eq:D.40-O}
\end{equation}
For the current, $s_\uparrow = s_\downarrow = 1$, and the bracket is $\Psi^\dagger_{\vect{k}}\Psi_{\vect{k}}$, the unit matrix in Nambu space. For the $z$ component of the spin current, $s_\sigma = \pm\frac12$, and the bracket is $\frac12\Psi^\dagger_{\vect{k}}\boldsymbol{\eta}\Psi_{\vect{k}}$, the matrix of the density. For an even weight the two assignments are exchanged: the density has $\boldsymbol{\eta}$ (\cref{app:D.21}), the spin density the unit matrix (\cref{sec:7.2}).

The inverse of \cref{eq:7.2-gamma} is $c_{\vect{k}\uparrow} = u\gamma_{\vect{k}1} + v\gamma^\dagger_{\vect{k}2}$ and $c^\dagger_{-\vect{k}\downarrow} = -v\gamma_{\vect{k}1} + u\gamma^\dagger_{\vect{k}2}$, with $u = u^{\mathrm{B}}_{\vect{k}}$ and $v = v^{\mathrm{B}}_{\vect{k}}$. Inserting it and using $\gamma_{\vect{k}2}\gamma^\dagger_{\vect{k}2} = 1 - \gamma^\dagger_{\vect{k}2}\gamma_{\vect{k}2}$, $u^2 + v^2 = 1$, $u^2 - v^2 = \xi_{\vect{k}}/E_{\vect{k}}$ and $2uv = \Delta_{\mathrm{BCS}}/E_{\vect{k}}$,
\begin{equation}
  \begin{aligned}
  \Psi^\dagger_{\vect{k}}\Psi_{\vect{k}} &= \gamma^\dagger_{\vect{k}1}\gamma_{\vect{k}1} - \gamma^\dagger_{\vect{k}2}\gamma_{\vect{k}2} + 1 , \\
  \Psi^\dagger_{\vect{k}}\boldsymbol{\eta}\Psi_{\vect{k}} &= \frac{\xi_{\vect{k}}}{E_{\vect{k}}}\bigl(\gamma^\dagger_{\vect{k}1}\gamma_{\vect{k}1} + \gamma^\dagger_{\vect{k}2}\gamma_{\vect{k}2} - 1\bigr) + \frac{\Delta_{\mathrm{BCS}}}{E_{\vect{k}}}\bigl(\gamma^\dagger_{\vect{k}1}\gamma^\dagger_{\vect{k}2} + \gamma_{\vect{k}2}\gamma_{\vect{k}1}\bigr) .
  \end{aligned}
  \label{eq:D.40-bogo}
\end{equation}
The current $\hat{\vect{e}}\cdot\op{\vect{V}} = \sum_{\vect{k}}(\hat{\vect{e}}\cdot\vect{v}_{\vect{k}})\Psi^\dagger_{\vect{k}}\Psi_{\vect{k}}$ is therefore diagonal in the quasiparticles and commutes with $H_{\mathrm{MF}}$. The quasiparticle $\gamma^\dagger_{\vect{k}1}$, an excitation of $(\vect{k}\!\uparrow)$, carries $\hat{\vect{e}}\cdot\vect{v}_{\vect{k}}$, and $\gamma^\dagger_{\vect{k}2}$, an excitation of $(-\vect{k}\!\downarrow)$, carries $-\hat{\vect{e}}\cdot\vect{v}_{\vect{k}} = \hat{\vect{e}}\cdot\vect{v}_{-\vect{k}}$; the ground state carries $\sum_{\vect{k}}\hat{\vect{e}}\cdot\vect{v}_{\vect{k}} = 0$. The spin current is not diagonal. Its last term creates the pair $\gamma^\dagger_{\vect{k}1}\gamma^\dagger_{\vect{k}2}|\mathrm{BCS}\rangle$ of energy $2E_{\vect{k}}$ and total momentum zero.

\subsection*{Spin current without impurities}

With $\hat{\vect{e}} = \hat{\vect{z}}$, the spin current $\mathcal{J}^{zz}_e = \frac12\sum_{\vect{k}}v^z_{\vect{k}}\Psi^\dagger_{\vect{k}}\boldsymbol{\eta}\Psi_{\vect{k}}$ connects $|\mathrm{BCS}\rangle$ to the pair state at $\vect{k}$ with the amplitude $\frac12v^z_{\vect{k}}\Delta_{\mathrm{BCS}}/E_{\vect{k}}$. Without spin--orbit coupling the Hamiltonian and the singlet ground state are invariant under rotations of the spin alone, so the three components $\mathcal{J}^{\alpha\alpha}_e$ of $\mathcal{J}_e$ do not interfere and each has the response of $\mathcal{J}^{zz}_e$ in a cubic crystal: $S_{\mathcal{J}\mathcal{J}} = 3S_{\mathcal{J}^{zz}\mathcal{J}^{zz}}$. The conductivity $\sigma_1^{\mathcal{J}} = (2e^2/3\omega)S_{\mathcal{J}\mathcal{J}}$ of \cref{eq:15.4-sJdef} is then
\begin{equation*}
  \sigma_1^{\mathcal{J}}(\omega) = \frac{2e^2}{\omega}\,\frac{2\pi}{\Vol}\sum_{\vect{k}}\frac14(v^z_{\vect{k}})^2\frac{\Delta_{\mathrm{BCS}}^2}{E_{\vect{k}}^2}\,\delta(\omega - 2E_{\vect{k}})
  = \frac{\pi e^2}{\omega}\,\mathcal{N}_{\mathrm{F}}\langle(v^z)^2\rangle_{\mathrm{FS}}\int\dd\xi\,\frac{\Delta_{\mathrm{BCS}}^2}{E^2}\,\delta(\omega - 2E) ,
\end{equation*}
where the sum over $\vect{k}$ has become $\Vol\mathcal{N}_{\mathrm{F}}\int\dd\xi$ times the Fermi-surface average, each $\vect{k}$ labeling one pair state of both spins. The delta function fixes $E = \omega/2$ and is reached from $\xi = \pm(\omega^2/4 - \Delta_{\mathrm{BCS}}^2)^{1/2}$, each with the weight $E/2|\xi|$. With the intraband plasma frequency $(\omega_p^{\mathrm{intra}})^2 = 2e^2\mathcal{N}_{\mathrm{F}}\langle(v^z)^2\rangle_{\mathrm{FS}}$ this gives \cref{eq:15.4-sJ}. Its integral follows with $\omega = 2\Delta_{\mathrm{BCS}}\cosh t$: $\dd\omega = 2\Delta_{\mathrm{BCS}}\sinh t\,\dd t$ and $(\omega^2 - 4\Delta_{\mathrm{BCS}}^2)^{1/2} = 2\Delta_{\mathrm{BCS}}\sinh t$, so that $\int_{2\Delta}^{\omega}\sigma_1^{\mathcal{J}}\dd\omega' = \frac{\pi}{2}(\omega_p^{\mathrm{intra}})^2\int_0^{t}\dd t'/\cosh^2t' = \frac{\pi}{2}(\omega_p^{\mathrm{intra}})^2\tanh t$. The total is $\frac{\pi}{2}(\omega_p^{\mathrm{intra}})^2$, the intraband f-sum weight.

The total weight, summed over all frequencies, does not depend on the mean-field treatment. For a Hamiltonian $H$ that does not act on spin, the $z$ spin current is the time derivative of the spin dipole $D_s = \sum_i\sigma^z_iz_i/2$, $\mathcal{J}^{zz}_e = \ii[H, D_s]$, and the double commutator $[[H, D_s], D_s]$ is a quarter of the one-body constant of the charge dipole, the factor that the normalization of \cref{eq:15.4-sJdef} compensates. The spin-current absorption therefore obeys the f-sum rule of \cref{sec:4.3} exactly, with interactions. Since $D_s$ commutes with every spin-independent term, this holds although $\mathcal{J}^{zz}_e$ itself is not conserved by $H$. The mean-field Hamiltonian respects it too: a contact singlet pairing term $\psi^\dagger_\uparrow\psi^\dagger_\downarrow$ commutes with the local spin density, so $H_{\mathrm{MF}}$ conserves spin locally, and the $s$-wave ladder of \cref{eq:7.2-Rc} does not dress the vertex $v^z_{\vect{k}}$, which is odd in $\vect{k}$. At $T = 0$ the paramagnetic part of the static response, $(2/\pi)\int\sigma_1^{\mathcal{J}}\dd\omega = (\omega_p^{\mathrm{intra}})^2$, cancels the diamagnetic term exactly, so the singlet carries no spin supercurrent, the counterpart of its vanishing Pauli susceptibility.

The weight near $2\Delta_{\mathrm{BCS}}$ is the coherent part of the total. Two corrections set it. The self-energy gives the quasiparticle residue $1/(1 + \lambda_{\mathrm{ep}})$ of \cref{eq:6.4-mass} (with $F_1^s$ entering through the mass); the phonon interaction does not depend on spin and adds nothing in the spin channel. The interaction in the $\ell = 1$ spin channel, the Landau parameter $F_1^a$, gives the quasiparticle spin current $(1 + F_1^a/3)\vect{v}^*$. Near the edge the weight is therefore $\frac{\pi}{2}(\omega_p^{\mathrm{intra}})^2(1 + F_1^a/3)/(1 + \lambda_{\mathrm{ep}})$. With $f(\omega)$ the bare bubble of \cref{eq:15.4-sJ} normalized to $f(0) = 1$, $\Imag f = 2\pi\Delta_{\mathrm{BCS}}^2/[\omega(\omega^2 - 4\Delta_{\mathrm{BCS}}^2)^{1/2}]$, the dressed kernel has the form of Leggett,
\begin{equation*}
  K_s(\omega) = \frac{(\omega_p^{\mathrm{intra}})^2}{1 + \lambda_{\mathrm{ep}}}\,\frac{(1 + F)\,[1 - f(\omega)]}{1 + F f(\omega)} , \qquad F = \frac{F_1^a}{3} , \qquad \sigma_1^{\mathcal{J}} = -\frac{\Imag K_s}{\omega} ,
\end{equation*}
which vanishes at $\omega = 0$ (no spin stiffness) and tends to the normal-state spin Drude weight for $f \to 0$. At $4\Delta_{\mathrm{BCS}}$ it multiplies the bare $\sigma_1^{\mathcal{J}}$ by 1.31, 1.15, 0.87 and 0.75 for $F = 0.1$, 0.05, $-0.05$ and $-0.1$. Below the edge $\Real f$ grows as $(2\Delta_{\mathrm{BCS}} - \omega)^{-1/2}$, so for $F_1^a < 0$ the denominator has a root below $2\Delta_{\mathrm{BCS}}$: a bound spin-one mode with $\ell = 1$, bound by $5.6\times10^{-3}$, $2.1\times10^{-2}$ and $7.4\times10^{-2}\Delta_{\mathrm{BCS}}$ at $F = -0.05$, $-0.1$ and $-0.2$, which takes weight from the continuum.

\subsection*{The phonon-assisted amplitude}

Let the band at the Fermi level couple to the phonons through \cref{eq:A.6-Hep} with the vertex $g_\nu(\vect{k},\vect{Q})$, band indices dropped, and assume the phase convention in which time reversal gives $g_\nu(-\vect{k}-\vect{Q},\vect{Q}) = g_\nu(\vect{k},\vect{Q})$, as for any vertex of a perturbation even under time reversal. The final states with one phonon emitted are
\begin{equation*}
  |f\rangle = \gamma^\dagger_{\vect{k}+\vect{Q},1}\,\gamma^\dagger_{\vect{k}2}\,b^\dagger_{-\vect{Q}\nu}|\mathrm{BCS}\rangle , \qquad E_f - E_i = E_{\vect{k}+\vect{Q}} + E_{\vect{k}} + \omega_{\vect{Q}\nu} ,
\end{equation*}
an excitation of $(\vect{k}+\vect{Q}\!\uparrow)$, one of $(-\vect{k}\!\downarrow)$ and the phonon; the other spin assignment is the same family with $\vect{k} \to -\vect{k}-\vect{Q}$. The phonon vertex creates them from $|\mathrm{BCS}\rangle$ through its spin-up term $c^\dagger_{\vect{k}+\vect{Q}\uparrow}c_{\vect{k}\uparrow}$, with the coefficient $u^{\mathrm{B}}_{\vect{k}+\vect{Q}}v^{\mathrm{B}}_{\vect{k}}$, and through its spin-down term $c^\dagger_{-\vect{k}\downarrow}c_{-\vect{k}-\vect{Q}\downarrow}$, with $v^{\mathrm{B}}_{\vect{k}+\vect{Q}}u^{\mathrm{B}}_{\vect{k}}$, exactly as the density does in \cref{eq:7.2-coh}; the time-reversal condition makes the two vertices equal, so that
\begin{equation*}
  \langle f|H_{\mathrm{ep}}|\mathrm{BCS}\rangle = \frac{1}{\sqrt N}\,g_\nu(\vect{k},\vect{Q})\,\bigl(u^{\mathrm{B}}_{\vect{k}+\vect{Q}}v^{\mathrm{B}}_{\vect{k}} + v^{\mathrm{B}}_{\vect{k}+\vect{Q}}u^{\mathrm{B}}_{\vect{k}}\bigr) .
\end{equation*}

The amplitude to reach $|f\rangle$ from $|i\rangle = |\mathrm{BCS}\rangle$ with one current vertex $O = \hat{\vect{e}}\cdot\op{\vect{V}}$ and one phonon vertex is \cref{eq:2.4-c2} with the two orderings of \cref{eq:13.2-T},
\begin{equation*}
  T_{fi} = \sum_m\left[\frac{\langle f|H_{\mathrm{ep}}|m\rangle\langle m|O|i\rangle}{E_i + \omega - E_m} + \frac{\langle f|O|m\rangle\langle m|H_{\mathrm{ep}}|i\rangle}{E_i - E_m}\right] , \qquad \omega = E_f - E_i ,
\end{equation*}
with $|m\rangle$ the eigenstates of $H_{\mathrm{MF}}$ plus the free phonons. Since $O$ is diagonal in that basis, $\langle m|O|i\rangle = J_i\delta_{mi}$ and $\langle f|O|m\rangle = J_f\delta_{fm}$, with $J_i$ and $J_f$ the eigenvalues of $O$. The sums collapse:
\begin{equation}
  T_{fi} = \frac{J_i\,\langle f|H_{\mathrm{ep}}|i\rangle}{\omega} + \frac{J_f\,\langle f|H_{\mathrm{ep}}|i\rangle}{E_i - E_f} = -\frac{J_f - J_i}{\omega}\,\langle f|H_{\mathrm{ep}}|i\rangle .
  \label{eq:D.40-T}
\end{equation}
The only input is that $O$ commutes with $H_{\mathrm{MF}}$ and with the phonon Hamiltonian. By \cref{eq:D.40-bogo}, $J_i = 0$ and $J_f = \hat{\vect{e}}\cdot(\vect{v}_{\vect{k}+\vect{Q}} - \vect{v}_{\vect{k}})$, which is \cref{eq:15.4-T}.

\subsection*{The phonon-assisted conductivity}

The absorbed power of a field $\hat{\vect{e}}\mathcal{E}\cos\omega t$ coupled through \cref{eq:6.5-HA} gives the conductivity as in \cref{eq:6.5-eps2dir}, $\sigma_1 = (\pi e^2/\omega\Vol)\sum_f|T_{fi}|^2\delta(E_f - E_i - \omega)$, where the sum over $f$ runs over $\vect{k}$, $\vect{Q}$ and $\nu$. The squared coherence factor of $\langle f|H_{\mathrm{ep}}|i\rangle$ is the case-I factor of \cref{eq:7.2-cohE}, $\frac12[1 - (\xi\xi' - \Delta_{\mathrm{BCS}}^2)/EE']$, with primes at $\vect{k}+\vect{Q}$. In the Fermi-surface approximation $g_\nu$ and $\vect{v}$ depend on the positions of $\vect{k}$ and $\vect{k}+\vect{Q}$ on the Fermi surface and not on $\xi$ and $\xi'$ over the range $\omega$, so the sum over $\vect{k}$ and $\vect{Q}$ factorizes into a Fermi-surface sum, carrying $\delta(\xi_{\vect{k}})\delta(\xi_{\vect{k}+\vect{Q}})$, times $\int\dd\xi\int\dd\xi'$. The term odd in $\xi\xi'$ integrates to zero. Each of $\xi$ and $\xi'$ reaches the quasiparticle energy $E$ from both signs, with $\dd\xi = E\,\dd E/(E^2 - \Delta_{\mathrm{BCS}}^2)^{1/2}$, and at fixed pair energy $x = \omega - \omega_{\vect{Q}\nu}$
\begin{multline*}
  \int\dd\xi\int\dd\xi'\,\frac12\Bigl(1 + \frac{\Delta_{\mathrm{BCS}}^2}{EE'}\Bigr)\delta(x - E - E') \\
  = 2\int_{\Delta_{\mathrm{BCS}}}^{x - \Delta_{\mathrm{BCS}}}\dd E\,\frac{E(x - E) + \Delta_{\mathrm{BCS}}^2}{(E^2 - \Delta_{\mathrm{BCS}}^2)^{1/2}[(x - E)^2 - \Delta_{\mathrm{BCS}}^2]^{1/2}}
  = 2x\,\mathfrak{g}_{\mathrm{I}}(x/\Delta_{\mathrm{BCS}}) ,
\end{multline*}
by \cref{eq:7.2-pairspec}; in the normal state it is $2x$, the two spin orientations of an electron--hole pair of energy $x$. The Fermi-surface sum of $|g_\nu|^2[\hat{\vect{e}}\cdot(\vect{v}_{\vect{k}+\vect{Q}} - \vect{v}_{\vect{k}})]^2\delta(\xi)\delta(\xi')\delta(\omega' - \omega_{\vect{Q}\nu})$ is $2\langle(\hat{\vect{e}}\cdot\vect{v})^2\rangle_{\mathrm{FS}}\mathcal{N}_{\mathrm{F}}\Omega N^2\,\alpha^2F_{\mathrm{tr}}(\omega')$ by the definition \cref{eq:15.4-a2Ftr}. With $\Vol = N\Omega$,
\begin{equation*}
  \sigma_1^{\mathrm{S}}(\omega) = \frac{\pi e^2}{\omega^3}\,2\langle(\hat{\vect{e}}\cdot\vect{v})^2\rangle_{\mathrm{FS}}\mathcal{N}_{\mathrm{F}}\int\dd\omega'\,\alpha^2F_{\mathrm{tr}}(\omega')\,2(\omega - \omega')\,\mathfrak{g}_{\mathrm{I}}\Bigl(\frac{\omega - \omega'}{\Delta_{\mathrm{BCS}}}\Bigr) ,
\end{equation*}
which is \cref{eq:15.4-holstein}. The pair energy $\omega - \omega'$ must exceed $2\Delta_{\mathrm{BCS}}$, which bounds $\omega'$ by $\omega - 2\Delta_{\mathrm{BCS}}$.

For the Debye form $\alpha^2F_{\mathrm{tr}} = 2\lambda_{\mathrm{tr}}\omega'^4/\omega_{\mathrm{D}}^4$ below $\omega_{\mathrm{D}}$, which satisfies $\lambda_{\mathrm{tr}} = 2\int\alpha^2F_{\mathrm{tr}}/\omega'\,\dd\omega'$, the normal-state rate $\gamma_{\mathrm{H}}$ of \cref{eq:15.4-holstein} needs $\int_0^{\omega}(\omega - \omega')\omega'^4\dd\omega' = \omega^6/30$ for $\omega < \omega_{\mathrm{D}}$ and $\int_0^{\omega_{\mathrm{D}}}(\omega - \omega')\omega'^4\dd\omega' = \omega\omega_{\mathrm{D}}^5/5 - \omega_{\mathrm{D}}^6/6$ above, which give the two forms of $\gamma_{\mathrm{H}}$ in \cref{sec:15.4-holstein}. The ratio of \cref{eq:15.4-holstein} to its normal-state value is then the ratio of $\int_0^{\omega - 2\Delta}\dd\omega'\,\omega'^4(\omega - \omega')\mathfrak{g}_{\mathrm{I}}$ to $\omega^6/30$, the form of Ref.~\cite[Eq.~(17)]{hochberg2016absorption}, whose elliptic integral is $\mathfrak{g}_{\mathrm{I}} = \mathrm{E}\bigl((1 - 4\Delta_{\mathrm{BCS}}^2/x^2)^{1/2}\bigr)$, the Landen transform (a change of modulus that maps one elliptic integral onto another) of the form below.

\subsection*{The coherence integrals in closed form}

With $y = x/\Delta_{\mathrm{BCS}}$ and $\Delta_{\mathrm{BCS}} = 1$, write $E = y/2 + s$. Then $E(y - E) = y^2/4 - s^2$ and $(E^2 - 1)[(y - E)^2 - 1] = (a^2 - s^2)(c^2 - s^2)$ with $a = y/2 - 1$ and $c = y/2 + 1$, so that the integrals of \cref{eq:7.2-pairspec} are $(2/y)\int_0^a\dd s\,(y^2/4 \pm 1 - s^2)/[(a^2 - s^2)(c^2 - s^2)]^{1/2}$. The substitution $s = a\sin\phi$ and the modulus $k = a/c = (y - 2)/(y + 2)$ (a number, not a momentum) turn them into $(2/yc)\int_0^{\pi/2}\dd\phi\,(y^2/4 \pm 1 - a^2\sin^2\phi)/(1 - k^2\sin^2\phi)^{1/2}$. With $\int_0^{\pi/2}\sin^2\phi\,(1 - k^2\sin^2\phi)^{-1/2}\dd\phi = (\mathrm{K} - \mathrm{E})/k^2$, with the complete elliptic integrals $\mathrm{K}(k)$ and $\mathrm{E}(k)$, and $a^2/k^2 = c^2$, the bracket becomes $(y^2/4 \pm 1 - c^2)\mathrm{K} + c^2\mathrm{E}$, where $y^2/4 + 1 - c^2 = -y$ and $y^2/4 - 1 - c^2 = -(y + 2)$. Hence
\begin{equation}
  \mathfrak{g}_{\mathrm{I}}(y) = \Bigl(1 + \frac2y\Bigr)\mathrm{E}(k) - \frac{4}{y + 2}\,\mathrm{K}(k) , \qquad
  \mathfrak{g}_{\mathrm{II}}(y) = \Bigl(1 + \frac2y\Bigr)\mathrm{E}(k) - \frac4y\,\mathrm{K}(k) ,
  \label{eq:D.40-g}
\end{equation}
for $y \ge 2$. At $y = 2$, $k = 0$ and $\mathrm{E} = \mathrm{K} = \pi/2$, so that $\mathfrak{g}_{\mathrm{I}} = \pi/2$ and $\mathfrak{g}_{\mathrm{II}} = 0$; as $y \to \infty$ both tend to one. At $y = 4$ they are 1.211 and 0.672, the values of \cref{sec:7.2}, and both forms agree with direct quadrature of \cref{eq:7.2-pairspec} to $10^{-7}$.

\subsection*{Impurity-assisted absorption}

In a disordered normal metal the one-electron eigenstates $\alpha$ have energies $\xi_\alpha$ and no definite momentum; time reversal pairs each with a partner $\bar\alpha$ of the same energy. The paired state is built from $(\alpha\!\uparrow, \bar\alpha\!\downarrow)$ as \cref{eq:7.2-BCS} is built from $(\vect{k}\!\uparrow, -\vect{k}\!\downarrow)$, with the same gap for nonmagnetic disorder (\cref{sec:15.3}) and with $u^{\mathrm{B}}_\alpha$, $v^{\mathrm{B}}_\alpha$ from \cref{eq:7.2-E} at $\xi_\alpha$. The current $\sum_{\alpha\beta\sigma}V^{e}_{\alpha\beta}c^\dagger_{\alpha\sigma}c_{\beta\sigma}$, with $V^{e}_{\alpha\beta} = \langle\alpha|\hat{\vect{e}}\cdot\op{\vect{v}}|\beta\rangle$, is odd under time reversal, $V^{e}_{\bar\alpha\bar\beta} = -V^{e}_{\beta\alpha}$. The step of \cref{eq:D.40-O}, with the relabeling $\alpha \to \bar\alpha$ in place of $\vect{k} \to -\vect{k}$, turns its spin-down term into $\sum V^{e}_{\alpha\beta}c_{\bar\alpha\downarrow}c^\dagger_{\bar\beta\downarrow}$, and the operator into $\sum_{\alpha\beta}V^{e}_{\alpha\beta}\Psi^\dagger_\alpha\Psi_\beta$, again the unit matrix. Its pair-creation amplitude from the ground state into $\gamma^\dagger_{\alpha1}\gamma^\dagger_{\beta2}$ is then $V^{e}_{\alpha\beta}(u^{\mathrm{B}}_\alpha v^{\mathrm{B}}_\beta - v^{\mathrm{B}}_\alpha u^{\mathrm{B}}_\beta)$, the case-II combination of \cref{eq:7.2-coh}, with the squared form $\frac12[1 - (\xi_\alpha\xi_\beta + \Delta_{\mathrm{BCS}}^2)/E_\alpha E_\beta]$. A spin current, even under time reversal, carries $\boldsymbol{\eta}$ instead and the case-I combination. With the squared matrix elements constant over the few $\Delta_{\mathrm{BCS}}$ where the coherence factors matter (\cref{sec:15.4-MB}), the two probes give $\mathfrak{g}_{\mathrm{II}}$ and $\mathfrak{g}_{\mathrm{I}}$. The assumption fails when the mean free path approaches the size of a pair, $\gamma_{\mathrm{rel}} \sim \Delta_{\mathrm{BCS}}$, where the correlation of $\alpha$ and $\beta$ through their momenta survives; the clean limit is \cref{eq:D.40-bogo}.
